% page 532 of the facsimile (image p-531 of the scan)
% block 157.5 x 242.0 mm ; left margin 8.0 mm
\pgc{-12.700mm}{0.000mm}{
\l{\cel{3}524.}
\l{}
\l{\sou{shamoto}. (\sou{mineral}.)Argilo specala, ek silikato di alumino,}
\l{\cel{3}qua debas sua karakteri a la maniero disua molekul-agre-}
\l{\cel{3}geso, e de qua onu facas la kruzeli en qui onu submisas}
\l{\cel{3}la metali ed altra substanci a la kaloro de furnazi.-}
\l{}
\l{}
\l{shampunar. (trans.)\cel{2}Fricionar, por netigar lu, la hararo}
\l{\cel{3}e la kranio-epidermo per lociono ek alkoholo e saponario,}
\l{\cel{3}kun parfum-aquo. - DEF.}
\l{}
\l{\sou{shancelar}. (\sou{netrans}.,\cel{1}\sou{pro,}\cel{1}de). Vacilar sur sua bazo, mar-}
\l{\cel{3}chanto. - F.}
\l{}
\l{\sou{shapko}. Kapo-vesto armeala, uzata (olim) che la Poloniani,}
\l{\cel{3}e quan +weris la lancieri di Napoleono IIIesma (Francia).}
\l{}
\l{}
\l{\sou{sharado}. Enigmato qua konssistas ye divinar vorto quan onu}
\l{\cel{3}povas partigar e di qua singla parto formacas altra vorto,}
\l{\cel{3}segun defino di la parto e di la toto. - DEFIRS.}
\l{}
\l{\sou{sharko}. (\sou{zool.)}\cel{2}Grosa fisho marala, tre devorema, ek la ge-}
\l{\cel{3}nero \textquotedbl{}squali\textquotedbl{}. - L. squalus carcharias. - E.}
\l{}
\l{\sou{sharlatano}. I. (a) Aventuristo qua iras de urbo ad urbo,}
\l{\cel{3}frapante sur tamburego por atraktar la turbo, fanfaro-}
\l{\cel{3}nante pri sua risanigi miraklatra, operacante e vendante}
\l{\cel{3}sua drogi sen bubiko, en la ferii e sur la placi. (b) Me-}
\l{\cel{3}diko empirikala di qua la remediili esas, segun lua aserto,}
\l{\cel{3}nofaliiva e preske panacea. - II. La persono qua explo-}
\l{\cel{3}tas la kredemeso di publiko (+\sou{publico}) - qua esforcadas}
\l{\cel{3}divenar famoza, populara. - DEFIRS.}
\l{}
\l{\sou{sharpo}. I. Larja bendo de stofo, +werata zone cirkum la kor-}
\l{\cel{3}po, cirkum la torso, de la shultro dextra a la hancho si-}
\l{\cel{3}nistra, ofixigita ye la tayo per nodo, kun insigno di ula}
\l{\cel{3}ofici statala, municipala, urbala, e c. - II. Bandajo,}
\l{\cel{3}+werata bandoliere, o - subtenate da la kolo - falanta sur}
\l{\cel{3}la pektoro, por subtenar la avanbrakio malada. - DFIR.}
\l{}
\l{\sou{sharpio}. Amaso de fili quin onu detiris de linjo par-konsu-}
\l{\cel{3}mita, e quan onu uzas por la pansi. - DFR.}
\l{}
\l{\sou{sheiko}. Chefo di tribuo de Arabi. - DEFIRS.}
\l{}
\l{\sou{sheklo}. Masho di kateno, di qua la formo ne esas olta di la}
\l{\cel{3}mashi cetera, ed uzesas por skopo specala (fixigo di hoko,}
\l{\cel{3}di sitelo en dragilo, e c.) - DE.}
\l{}
\l{\sou{shelo.}\cel{2}I. (\sou{zool.}) Envelopilo kalka di la ovo. Envelopilo}
\l{\cel{3}kalka di ula moluski. - III. (\sou{bot.}) Envelopilo di ula}
\l{\cel{3}frukti (nuci, e c.). IV. (\sou{bot.)}\cel{2}Envelopilo quan onu de-}
\l{\cel{3}prenas de la frukti, de la legumi, kom preparo por lia}
\l{\cel{3}konsumeso. - DER.}
}
